\chapter{Introducción}
\orientacion{La introducción de la Especificación de requisitos de software (SRS) proporciona una vista general del documento e incluye su propósito, alcance, definiciones y acrónimos, referencias y organización.}

\section{Propósito}
\campo{Inserte aquí el texto}
\begin{itemize}
    \item Propósito del documento.
    \item Audiencia a la que va dirigido.
\end{itemize}

\section{Alcance}
\campo{Inserte aquí el texto}
\begin{itemize}
    \item Identificación del producto o de los productos que se desarrollarán mediante un nombre.
    \item Consistencia con definiciones similares de documentos de mayor nivel que puedan existir, por ejemplo, la descripción del sistema.
\end{itemize}

\section{Personal involucrado}
\begin{longtable}{p{3.4cm}p{10.2cm}}
\toprule
\textbf{Dato} & \textbf{Información} \\
\midrule
Nombre & \campo{Inserte aquí el texto} \\
Rol & \campo{Inserte aquí el texto} \\
Categoría profesional & \campo{Inserte aquí el texto} \\
Responsabilidades & \campo{Inserte aquí el texto} \\
Información de contacto & \campo{Inserte aquí el texto} \\
Aprobación & \campo{Inserte aquí el texto} \\
\bottomrule
\end{longtable}

\orientacion{Relación de las personas involucradas en el desarrollo del sistema, con información de contacto. Esta información permite localizar a los participantes y obtener la información necesaria para la obtención de requisitos y las validaciones de seguimiento.}

\section{Definiciones, acrónimos y abreviaturas}
\campo{Inserte aquí el texto}
\orientacion{Defina todos los términos, abreviaturas y acrónimos necesarios para interpretar apropiadamente este documento. Puede incluir referencias a uno o más apéndices u otros documentos.}

\section{Referencias}
\begin{longtable}{p{2.1cm}p{3.6cm}p{3.6cm}p{2.0cm}p{2.3cm}}
\toprule
\textbf{Referencia} & \textbf{Título} & \textbf{Ruta} & \textbf{Fecha} & \textbf{Autor} \\
\midrule
\campo{Referencia} & \campo{Título} & \campo{Ruta} & \campo{Fecha} & \campo{Autor} \\
\bottomrule
\end{longtable}

\orientacion{Relación completa de los documentos relacionados con la especificación de requisitos de software. Para cada uno, identifique el título, la referencia (si procede), la fecha y la organización que lo proporciona.}

\section{Resumen}
\campo{Inserte aquí el texto}
\orientacion{Describa el contenido del resto del documento y explique cómo está organizado.}
